% =================================================================================================
% APÊNDICES A E B: MANUAIS DE INSTALAÇÃO E DE USO DO DIAGNOSISXAI
% -------------------------------------------------------------------------------------------------
% Onde colar: no trabalho.tex, logo DEPOIS destas duas linhas e ANTES de \end{document}:
%
%     % Carrega o arquivo referencias.bib
%     \bibliography{referencias}
%
% Regras do manual da UTP (Normas Técnicas, 4. ed., 2024) aplicadas aqui:
%  - apêndice é texto elaborado pelo próprio autor; "os manuais de instrução ou
%    instalação do produto prototipado" são um dos exemplos do manual (p. 48);
%  - vem depois das referências (estrutura da monografia, p. 143);
%  - identificação por letra maiúscula consecutiva, travessão e título, tudo em
%    maiúsculas: "APÊNDICE A – QUESTIONÁRIO..." (p. 49);
%  - título sem indicativo numérico, centralizado (p. 16), começando em página
%    nova, como os de nível 1 (p. 16); o comando \apendice do template já faz isso;
%  - o título vai em MAIÚSCULAS também no argumento de \apendice, porque o
%    sumário repete o texto do argumento e deve ter a "mesma grafia" do texto
%    (p. 42); o exemplo de sumário do manual traz os apêndices em maiúsculas (p. 43);
%  - subtítulos de nível 2 em MAIÚSCULAS, sem negrito, indicativo separado do
%    título por um espaço, sem traço nem parêntese (p. 16);
%  - alíneas: texto anterior termina em dois-pontos; cada alínea começa com
%    letra minúscula e termina em ponto e vírgula; a última, em ponto final;
%  - itálico só para palavras estrangeiras (p. 20), como software, notebook e
%    Releases; aspas para destacar uma expressão (p. 19), aqui a mensagem
%    literal exibida pelo Windows.
%
% Conferido em 08/11/2026 contra o repositório Gabas/DiagnosisXAI (commit 22cbe5a),
% o LEIA-ME do empacotamento e o Release v1.0 (arquivos DiagnosisXAI-windows-x64.zip e
% DiagnosisXAI-linux-x64.zip). Um Release novo (v1.1, recomendado, ver
% 01_pendencias_email.md, item 18) mantém os mesmos nomes de arquivo.
% =================================================================================================

\apendice{MANUAL DE INSTALAÇÃO DO DIAGNOSISXAI}

Este apêndice descreve como instalar e executar o DiagnosisXAI, \textit{software}
desenvolvido neste trabalho. Há duas formas de uso: o executável, indicado a
quem deseja apenas operar o sistema, e o código-fonte, necessário para
reproduzir os experimentos ou retreinar os modelos. O código-fonte e os
executáveis estão disponíveis no repositório
\url{https://github.com/Gabas/DiagnosisXAI}.

\section*{A.1 REQUISITOS}

Para utilizar o sistema, são necessários:

\begin{alineas}
    \item computador com Windows ou Linux de 64 bits, para o executável, ou com
          Windows, Linux ou macOS, para o código-fonte;
    \item tela com resolução de 1920~$\times$~1080 pixels ou próxima dela; em
          telas menores, cada janela é limitada a 90\% da área disponível;
    \item navegador web, para as versões interativas do mapa populacional e da
          margem do SVM;
    \item o Python 3.12 e o Git, somente para a execução pelo código-fonte.
\end{alineas}

\section*{A.2 INSTALAÇÃO DO EXECUTÁVEL}

O executável dispensa a instalação do Python e de qualquer biblioteca. Para
instalá-lo:

\begin{alineas}
    \item acesse a página de versões (\textit{Releases}) do repositório,
          \url{https://github.com/Gabas/DiagnosisXAI/releases}, e baixe o
          arquivo DiagnosisXAI-windows-x64.zip ou DiagnosisXAI-linux-x64.zip
          da versão mais recente, conforme o sistema operacional;
    \item extraia o conteúdo do arquivo .zip em uma pasta de sua escolha,
          mantendo todos os arquivos juntos, pois o programa depende da pasta
          inteira, e não apenas do executável;
    \item no Windows, abra o arquivo DiagnosisXAI.exe; caso o sistema exiba o
          aviso ``O Windows protegeu o computador'', emitido porque o programa
          não tem assinatura digital, selecione Mais informações e, em seguida,
          Executar assim mesmo;
    \item no Linux, abra um terminal na pasta extraída e execute
          \texttt{./DiagnosisXAI}; se necessário, conceda antes a permissão de
          execução com \texttt{chmod +x DiagnosisXAI};
    \item opcionalmente, para conferir a instalação, execute o programa com a
          opção \texttt{-}\texttt{-autoteste}, que percorre sem interface as
          catorze etapas do uso real e grava o resultado no arquivo
          autoteste.txt, em uma pasta temporária.
\end{alineas}

A primeira abertura leva alguns segundos. A pasta exemplos, que acompanha o
executável, traz o conjunto de teste da base Wisconsin em dois arquivos:
dataTeste\_sem\_diagnostico.csv, para o primeiro passo, e
dataTeste\_com\_diagnostico.csv, gabarito para a auditoria do quarto passo.

O executável grava o histórico de sessões e os relatórios em PDF na pasta
DiagnosisXAI do diretório pessoal do usuário (C:\textbackslash
Users\textbackslash \textless usuário\textgreater \textbackslash DiagnosisXAI,
no Windows, e \textasciitilde/DiagnosisXAI, no Linux), pois a pasta do programa
pode ser somente leitura. Para desinstalar o sistema, basta apagar a pasta
extraída e, se desejado, essa pasta de dados.

\section*{A.3 EXECUÇÃO A PARTIR DO CÓDIGO-FONTE}

A execução pelo código-fonte é indicada para reproduzir os experimentos do
Capítulo~\ref{cap:metodologia}. No macOS e no Linux, a biblioteca gráfica
Tkinter, usada pela interface, é distribuída separadamente do Python e deve ser
instalada junto com ele, por exemplo pelo pacote python3-tk, no Ubuntu, ou
python-tk@3.12, no Homebrew; no Windows, o instalador oficial do Python já a
inclui. Os passos são:

\begin{alineas}
    \item clone o repositório com o comando \texttt{git clone}, seguido do
          endereço \url{https://github.com/Gabas/DiagnosisXAI.git}, e entre na
          pasta DiagnosisXAI;
    \item crie um ambiente virtual com \texttt{python3 -m venv .venv} e
          ative-o com \texttt{source .venv/bin/activate}, no Linux e no macOS,
          ou \texttt{.venv\textbackslash Scripts\textbackslash activate}, no
          Windows;
    \item instale as dependências com
          \texttt{pip install -r requirements.txt}; para reproduzir exatamente o
          ambiente que gerou o arquivo wisconsin.pkl, use o arquivo
          requirements-executavel.txt, que fixa as versões exatas;
    \item entre na pasta app e inicie a aplicação com
          \texttt{python main.py}.
\end{alineas}

Os modelos treinados já acompanham o repositório, no arquivo
data/wisconsin.pkl. Para retreiná-los, abra o \textit{notebook}
notebooks/Wisconsin.ipynb no Jupyter e execute todas as células, o que regrava o
wisconsin.pkl; em seguida, execute \texttt{python scripts/calibrar\_limiares.py}
a partir da raiz do repositório, para recalcular os limiares e as faixas de
recusa gravados em data/limiares.json. Os testes automatizados são executados
com o comando \texttt{pytest}, também a partir da raiz.

\apendice{MANUAL DE USO DO DIAGNOSISXAI}

Este apêndice orienta a operação do DiagnosisXAI. As telas mencionadas estão
ilustradas na Seção~\ref{sec:software}.

\section*{B.1 ORGANIZAÇÃO DA INTERFACE}

A barra lateral dá acesso a quatro abas:

\begin{alineas}
    \item a aba Início exibe o estado do sistema, como os modelos carregados;
    \item a aba Novo Diagnóstico conduz o processamento de um lote em seis
          passos;
    \item a aba Histórico lista as sessões anteriores e permite reabrir seus
          relatórios de explicabilidade sem reprocessar o lote;
    \item a aba Sobre reúne o memorial de cálculo, que explica a origem de cada
          número exibido, e o glossário dos trinta biomarcadores.
\end{alineas}

O tema claro ou escuro da interface pode ser trocado no seletor Tema, no rodapé
da barra lateral.

\section*{B.2 PREPARAÇÃO DA PLANILHA DE ENTRADA}

O sistema recebe os pacientes em um arquivo CSV, uma linha por paciente, com as
trinta colunas morfológicas da base Wisconsin, nomeadas como na base original
(radius\_mean, texture\_mean, \ldots, fractal\_dimension\_worst). Os valores
devem ser os brutos, medidos pelo laboratório, e não padronizados, pois a
padronização é feita pelo próprio sistema. A planilha não deve conter as colunas
id nem diagnosis. O lote é recusado, com mensagem que indica a coluna e a linha
do problema, quando falta alguma coluna, há células vazias, texto no lugar de
número, valores negativos ou nenhuma linha de dados.

Para a auditoria do quarto passo, usa-se um segundo arquivo, o gabarito, com as
mesmas colunas acrescidas da coluna diagnosis, preenchida com 1 para maligno e
0 para benigno, como no arquivo dataTeste\_com\_diagnostico.csv.

\section*{B.3 PROCESSAMENTO DE UM LOTE}

Na aba Novo Diagnóstico, cada passo é liberado somente após a conclusão do
anterior, conforme a Figura~\ref{fig:app_pipeline}:

\begin{alineas}
    \item no Passo 1, Importar Dados Brutos (CSV), selecione a planilha; os
          valores aparecem na tabela;
    \item no Passo 2, Higienizar e Escalar (Z-Score), clique em Aplicar
          Padronização; o sistema valida a planilha e aplica a padronização
          ajustada no conjunto de treino;
    \item no Passo 3, Inteligência Artificial, escolha na lista um dos cinco
          classificadores, o Comitê (voto suave) ou Todos (Comparação), que é a
          opção inicial; confira a régua de decisão exibida e clique em
          Processar Diagnóstico; a opção Adiar casos incertos em vez de decidir
          vem ligada e só está disponível para a Regressão Logística, o Random
          Forest, o SVM e o comitê;
    \item no Passo 4, Auditoria Acadêmica, que é opcional, carregue o gabarito
          para obter a cobertura, as métricas de desempenho, a matriz de
          confusão e uma leitura crítica dos resultados;
    \item no Passo 5, Explicabilidade (XAI), abra a janela de explicação do
          modelo executado e selecione um paciente para ver os fatores que mais
          pesaram e o raciocínio próprio do modelo; nessa janela, o botão
          Exportar Paciente (PDF) gera o laudo individual do paciente;
    \item no Passo 6, Exportar Resultados, salve a tabela completa em CSV ou o
          relatório do lote em PDF, acompanhado da régua de decisão.
\end{alineas}

\section*{B.4 LEITURA DOS RESULTADOS}

A tabela de resultados traz, para cada paciente, as seguintes colunas:

\begin{alineas}
    \item a coluna Diagnóstico\_IA indica Maligno, Benigno ou Revisar; Revisar
          não é uma classe intermediária, e sim a indicação de que o sistema se
          absteve e de que o caso deve ser avaliado pelo especialista;
    \item a coluna Certeza\_Maligno(\%) mostra a probabilidade calibrada de
          malignidade, que não deve ser comparada a 50\%, pois o diagnóstico
          decorre dos cortes da régua exibida no terceiro passo;
    \item a coluna Zona\_de\_Decisão indica Definida, Revisar ou, quando a
          recusa está desligada, Limítrofe, para os casos a menos de dez pontos
          percentuais do limiar de operação;
    \item a coluna Perfil indica Típico ou Atípico; Atípico significa que o
          paciente se afasta da população de treino, situação em que a
          predição merece cautela, e não que o tumor seja maligno.
\end{alineas}

No modo Todos (Comparação), essas colunas são substituídas por uma coluna por
modelo, com o diagnóstico que cada um daria ao mesmo paciente.

\section*{B.5 CUIDADOS DE USO}

O DiagnosisXAI é um \textit{software} acadêmico de apoio, não é um dispositivo
médico e não substitui o diagnóstico do profissional de saúde. O sistema não
processa imagens: opera apenas sobre os trinta atributos já extraídos pelo
laboratório, obtidos pelo mesmo procedimento da base Wisconsin. Os casos
devolvidos para revisão e os pacientes de perfil atípico devem ser examinados
com atenção redobrada.
